\documentstyle[%


righttag%


,11pt%


,amssymb%


,russian%               remove in Eng.


,izvru%                 Set izven% in Eng.


% ,izvru-ks             for brief communications


]{amsart}





%       Math definitions





\newcommand{\la}{\langle}


\newcommand{\ra}{\rangle}


% \newcommand{\wt}{\widetilde}


\newcommand{\Cl}{\operatorname{Cl}}


\newcommand{\card}{\operatorname{card}}


\newcommand{\supp}{\operatorname{supp}}


\newcommand{\vraisup}{\operatorname{vrai~sup}}


\newcommand{\tts}[1]{}





\theoremstyle{plain}


\theorembodyfont{\it}


\newtheorem{thms}{\hskip\parindent Теорема}[section]





\theoremstyle{definition}


\theorembodyfont{\rm}


\newtheorem{examnonum}{\hskip\parindent Пример}


\newtheorem{cors}{\hskip\parindent Следствие}[section]


\renewcommand{\theexamnonum}{}





%\hoffset=-15mm                  For Eng.





\setcounter{page}{54}


\god{1998}


\nomer{2~(429)} %                Remove in Eng.


%\nomer{Vol.\,42, No.\,2}        Set in Eng.


% \str{75--81}                    Set in Eng.


\udk{517.982}


\author{\it А.Г.\,Мясников}


% NB    ^^ remove in Eng.


\title{Аменабельные $L_1(G)$-модули, усредняемые функции\\


и пространства со смешанной нормой}





\thanks{Работа поддержана Государственным комитетом Российской Федерации


по высшему образованию (грант \No\,216424).}


\date{}


% \pagestyle{myheadings}         Set in Eng.





\pagestyle{plain} %              Remove in Eng.


\begin{document}





% \thispagestyle{plain}          Set in Eng.





\maketitle





\section{Введение}





В настоящей статье продолжается исследование аменабельных модулей [1].


Основное внимание уделяется изучению свойств топологически усредняемых


функций из компоненты аменабельности в $L_\infty(G)$. Полученные


результаты


применяются для решения вопроса об аменабельности функциональных модулей


и, в частности, модулей со смешанными $L_p$-нормами. Дополнительная


информация об усредняемых функциях может быть найдена в работах


[2]--[4].





В дальнейшем через $G$ обозначается локально компактная группа с левой


мерой Хаара $\lambda_G$, через $L_p(G)$, $p\in [1,\infty]$, ---


соответствующие пространства


комплекснозначных функций, рассматриваемые как левые банаховы


$L_1(G)$-модули


относительно свертки


$\varphi\ast f=\int\limits_G\varphi(s)f(s^{-1}t)d\lambda_G(s)$,


$t\in G$ ($\varphi\in L_1(G)$,


$f\in L_p(G)$). Далее, $P(G)=\{\varphi\in L_1(G):\varphi\ge 0,\ \;


\|\varphi\|_1=1\}$, $RUC(G)=L_1(G)\ast L_\infty(G)$, $C_{00}(G)$ ---


множество


непрерывных на $G$ комплекснозначных функций с компактными носителями.


Значение функционала $F\in X^*$ в $x\in X$ обозначается


$\la x,F\ra$. И, наконец,


используются обозначения


${}_sf(t)=f(st)$, $f_s(t)=f(ts)$, $\wt f(t)=f(t^{-1})$,


$f^*(t)=\Delta_G(t^{-1})f(t^{-1})$, где $\Delta_G$ ---


модулярная функция.





\section{Топологически левоусредняемые функции}





Левый банахов $L_1(G)$-модуль $M$


называется аменабельным, если множество всех


$x\in M$, удовлетворяющих условию


$\inf\{\|\varphi x\|:\varphi\in P(G)\}=0$,


замкнуто относительно сложения (свойство


Эмерсона [5]). Произвольный левый банахов $L_1(G)$-модуль $M$ содержит


наибольший аменабельный подмодуль $AM$, который называется компонентой


аменабельности в $M$. Наконец, модуль $\cal L_M(G)$ определяется как


замкнутый


подмодуль в $L_\infty(G)$, порожденный функциями


$x\circ F$, $\overline{x\circ F}$ ($x\in M$, $F\in M^*$), где


функция $x\circ F$ определена равенством


$\la\varphi,x\circ F\ra=\la\varphi^*x,F\ra$ ($\varphi\in L_1(G)$).





Значение последнего определения заключается в том, что аменабельность


$M$


эквивалентна аменабельности $\cal L_M(G)$, а также эквивалентна


существованию


топологически левоинвариантного на $\cal L_M(G)$ среднего (функционал


$\Phi\in L_\infty(G)^*$ называется


топологически левоинвариантным на $X\subset L_\infty(G)$ средним, если


1)~$\Phi\ge 0$; 2)~$\|\Phi\|=1$; 3)~$\la\varphi\ast f,\Phi\ra=\la


f,\Phi\ra$ для всех $f\in X$, $\varphi\in P(G)$) [1].





Функция $f\in L_\infty(G)$ называется топологически левоусредняемой к


постоянной $c$, если замыкание множества


$P(G)\ast f=\{\varphi\ast f:\varphi\in P(G)\}$ содержит функцию


$c1_G$ ($1_G(t)\equiv t$).





Пусть $f\in L_\infty(G)$. Обозначим через $T_f$ замкнутое линейное


подпространство в $L_\infty(G)$,


порожденное функциями $f$ и $\varphi\ast f$ ($\varphi\in L_1(G)$).


Будем говорить, что функция $f$ имеет


единственное топологически левоинвариантное среднее значение, если


множество топологически левоинвариантных на $T_f$ средних не пусто,


причем все такие средние принимают в $f$ одно и то же значение.





\begin{thms}


Пусть $f\in L_\infty(G)$, тогда следующие условия эквивалентны{\em :}


\begin{itemize}


\item[1)] функция $f$ принадлежит $AL_\infty(G)$ и топологически


левоусредняема к постоянной $c;$


\item[2)] для каждого $\varphi\in P(G)$ свертка $\varphi\ast f$


топологически


левоусредняема к единственной постоянной $c$ {\em (}не зависящей от


$\varphi);$


\item[3)] функция $f$ принадлежит $AL_\infty(G)$ и имеет единственное


топологически левоинвариантное среднее значение $c;$


\item[4)] для каждого среднего $F\in L_\infty(G)^*$ функция


$f\circ F$ имеет


единственное


топологически левоинвариантное среднее значение $c$


{\em (}не зависящее от $F)$.


\end{itemize}


\end{thms}





\begin{pf}


Эквивалентность условий 1)--4) достаточно установить для случая $c=0$.





$2)\,\Longrightarrow\,1)$, $1)\,\Longrightarrow\,4)$. Очевидно.





$4)\,\Longrightarrow\,3)$. Каждому конечному подмножеству


$X\subset\cal L_{T_f}(G)$ и $\varepsilon>0$


сопоставим среднее $\Phi_{X,\varepsilon}$ такое,


что $|\{\Phi_{X,\varepsilon}(g)\}|<\varepsilon$ для всех $g\in X$. Пусть


$\Phi$


--- предельная точка сети $\{\Phi_{X,\varepsilon}\}$ в $w^*$-топологии.


Тогда $\Phi(g)=0$ для всех $g\in\cal L_{T_f}(G)$, следовательно, модуль


$T_f$ аменабелен и, в частности, $f\in AL_\infty(G)$.





Пусть $\{\varphi_\alpha\}\subset P(G)$ --- аппроксимативная единица для


$L_1(G)$, $I\in L_\infty(G)^*$ --- предельная


точка сети $\{\varphi_\alpha\}$ в $w^*$-топологии. Тогда $f=f\circ I$,


следовательно, функция $f$


имеет единственное топологически левоинвариантное среднее значение $0$.





$3)\,\Longrightarrow\,2)$. Фиксируем среднее $\Phi$, а также


топологически левоинвариантное на $AL_\infty(G)$


среднее $\Psi$. Определим топологически левоинвариантное на


$AL_\infty(G)$ среднее $\Phi\circ\Psi$


равенством


$\la g,\Phi\circ\Psi\ra=\la g\circ\Phi,\Psi\ra$ ($g\in L_\infty(G)$).


Пусть $\Psi=w^*-\lim\limits_\beta\psi_\beta$, где $\psi_\beta\in P(G)$.


Тогда для любого $\varphi\in P(G)$


$$


\lim_\beta\la\psi^*_\beta\ast\varphi\ast f,\Phi\ra=


\la\varphi\ast f,\Phi\circ\Psi\ra=


\la f,\Phi\circ\Psi\ra=0,


$$


откуда $w-\lim\limits_\beta\psi^*_\beta\ast\varphi\ast f=0$. Далее


рассуждаем, как в ([6], сс.\,48, 49).


\end{pf}





Подмодуль $M\subset L_\infty(G)$ называется топологически


левоинтровертным, если $x\circ F\in M$ для всех


$x\in M$, $F\in M^*$. Очевидно, множество всех функций, удовлетворяющих


условиям 1)--4)


теоремы 2.1, является топологически левоинтровертным подмодулем в


$AL_\infty(G)$.


Заметим также, что в случае $c=0$ требование единственности в условии 2)


выполняется автоматически и может быть опущено.





\begin{examnonum}


Любая слабо почти периодическая функция на $G$ принадлежит


$AL_\infty(G)$ ([1], пример 2.3) и имеет единственное топологически


левоинвариантное среднее значение


([7], с.86), следовательно, удовлетворяет условиям теоремы 2.1. Пусть


теперь модуль $M$ таков, что для каждого $x\in M$ множество $P(G)x$


является


относительно компактным в $w$-топологии. Тогда для любых


$x\in M$, $F\in M^*$ множество


функций $P(G)\ast(x\circ F)$ также относительно компактно в


$w$-топологии и, следовательно,


состоит из слабо почти периодических функций ([8], теорема~4). Таким


образом, функции из $\cal L_M(G)$ удовлетворяют условиям теоремы 2.1.


\end{examnonum}





\begin{cors}


Для произвольного левого банахова $L_1(G)$-модуля $M$ следующие условия


эквивалентны:


\begin{itemize}


\item[1)] найдется сеть $\{\varphi_\alpha\}\subset P(G)$ такая, что


$\lim\limits_\alpha\|\varphi_\alpha x\|=0$, $x\in M$;


\item[2)] каждая функция из $\cal L_M(G)$ топологически левоусредняема


к нулю.


\end{itemize}


\end{cors}





\begin{pf}


$1)\,\Longrightarrow\,2)$. Очевидно.





$2)\,\Longrightarrow\,1)$. По теореме 2.1 модуль $\cal L_M(G)$


аменабелен, следовательно, найдется среднее


$\Phi$ такое, что $\Phi(f)=0$, $f\in\cal L_M(G)$. Пусть


$\Phi=w^*-\lim\limits_\beta\varphi_\beta$, $\varphi_\beta\in P(G)$.


Тогда $w-\lim\limits_\beta\varphi^*_\beta x=0$ для всех $x\in M$.


Рассуждая, как в ([6], сс.\,48, 49), убедимся в выполнении условия 1).


\end{pf}





\section{Ядро компоненты аменабельности и функциональные модули}





Назовем ядром компоненты аменабельности в $L_\infty(G)$ множество


$\cal K(G)$ функций $f\in L_\infty(G)$


таких, что свертка $\varphi\ast |f|$ топологически левоусредняема к


нулю для любого


$\varphi\in P(G)$. Из теоремы 2.1 следует, что $\cal K(G)$ является


замкнутым подмодулем в $AL_\infty(G)$.





\begin{thms}


Пусть группа $G$ не является аменабельной и $f\in L_\infty(G)$. Тогда


функция


$f$ принадлежит $\cal K(G)$ в том и только том случае, когда любая


функция $g\in L_\infty(G)$, удовлетворяющая условию $0\le g\le |f|$,


принадлежит $AL_\infty(G)$.


\end{thms}





\begin{pf}


Пусть $f\in\cal K(G)$ и $0\le g\le |f|$. Так как


$\varphi\ast g\le\varphi\ast |f|$, $\varphi\in P(G)$, то функция $g$


удовлетворяет условиям теоремы 2.1. В частности, $g\in AL_\infty(G)$.





Предположим теперь, что любая функция $g\in L_\infty(G)$,


удовлетворяющая условию


$0\le g\le |f|$, принадлежит $AL_\infty(G)$. Очевидно,


$|f|\in AL_\infty(G)$. Предположим,


что $\Phi(|f|)>0$ для некоторого


топологически левоинвариантного на $AL_\infty(G)$ среднего $\Phi$.


Определим отображение


$\Phi_0:RUC(G)_+\to\Bbb R_+$ равенством


$\Phi_0(h)=\sup\{\Phi(g):g\in ARUC(G),\ \; 0\le g\le h\}$. Продолжая


$\Phi_0$ линейно на $RUC(G)$, получим ненулевой


топологически левоинвариантный на $RUC(G)$ функционал. Однако


существование


такого функционала противоречит предположению о неаменабельности $G$.


\end{pf}





\begin{thms}


Ядро $\cal K(G)$ является топологически левоинтровертным подмодулем в


$L_\infty(G)$.


\end{thms}





\begin{pf}


Пусть $f\in\cal K(G)$, $\Phi\in L_\infty(G)^*$ --- среднее и $\Psi$ ---


топологически левоинвариантное на $AL_\infty(G)$ среднее. Тогда


$|f|\circ\Phi\in AL_\infty(G)$ и т.к. $\Phi\circ\Psi$ --- топологически


левоинвариантное на $AL_\infty(G)$ среднее, то


$\la |f|\circ\Phi,\Psi\ra=\la |f|,\Phi\circ\Psi\ra=0$.


По теореме 2.1\; $|f|\circ\Phi\in\cal K(G)$. Из неравенства


$|f\circ\Phi|\le |f|\circ\Phi$ и теоремы


3.1 следует, что $f\circ\Phi\in\cal K(G)$.


\end{pf}





Покажем теперь, как ядро компоненты аменабельности в $L_\infty(G)$


может быть


использовано для выяснения вопроса об аменабельности функциональных


$G$-модулей. В оставшейся части статьи в целях простоты предполагается,


что группа $G$ является $\sigma$-компактной.





Обозначим через $\cal L(G)$ множество всех $\lambda_G$-измеримых


комплекснозначных функций.


При этом функции, совпадающие почти всюду, отождествляются. Отображение


$\rho:\cal L(G)\to[0,\infty]$ называется функциональной нормой на


$\cal L(G)$, если


выполнены следующие


условия: 1)~$\rho(f)=0$ в том и только том случае, когда $f=0$;


2)~$\rho(\alpha f)=|\alpha|\rho(f)$; 3)~$\rho(f+g)\le\rho(f)+\rho(g)$;


4)~$|f|\le |g|$


влечет $\rho(f)\le\rho(g)$. Функциональная норма $\rho$ определяет


нормированное линейное


пространство $L_\rho(G)=\{f:\rho(f)<\infty\}$, называемое нормированным


функциональным


пространством (или нормированным идеальным пространством).





Функциональная норма $\rho$ называется абсолютно непрерывной, если из


условия $f_k\downarrow 0$,\linebreak


$f_k\in L_\rho(G)$, $k=1,2,\dotsc,$ следует


$\rho(f_k)\downarrow 0$.





Определим ассоциированную норму


$\rho'(g)=\sup\Bigl\{\int\limits_G|fg|d\lambda_G:\rho(f)\le 1\Bigr\}$.


Если норма $\rho$ абсолютно


непрерывна, то нормируемое пространство $L_{\rho'}(G)$ является


сопряженным к $L_\rho(G)$ с двойственностью


$\la f,g\ra=\int\limits_G fg\,d\lambda_G$.





Пусть $\rho$ --- функциональная норма на $\cal L(G)$. Нормированное


пространство $L_\rho(G)$


будем называть банаховым функциональным $G$-модулем, если $L_\rho(G)$


полно, $\rho({}_sf)=\rho(f)$ и


$f_s\in L_\rho(G)$ для всех $f\in L_\rho(G)$, $s\in G$.


Отметим следующие свойства


банахова функционального $G$-модуля $L_\rho(G)$ с абсолютно непрерывной


нормой $\rho$:


\begin{itemize}


\item[1)] функция $\omega_\rho(s)=\sup\{\rho(f_s):\rho(f)\le 1\}$,


$s\in G$, полунепрерывна снизу и локально ограничена;


\item[2)] $\varphi\ast f\in L_\rho(G)$ и


$\rho(\varphi\ast f)\le \|\varphi\|_1\rho(f)$ для всех


$\varphi\in L_1(G)$, $f\in L_\rho(G)$;


\item[3)] $f\ast\varphi\in L_\rho(G)$ и


$\rho(f\ast\varphi)\le\|\wt\omega_\rho\wt\Delta_G\varphi\|_1


\rho(f)$ для всех


функций $\varphi\in L_1(G)$ с компактным носителем и всех


$f\in L_\rho(G)$;


\item[4)] множество $C_{00}(G)$ плотно в $L_\rho(G)$ ([9], теорема~3).


\end{itemize}





Примем обозначение $\cal L_\rho(G)=\cal L_{L_\rho(G)}(G)$.





\begin{thms}


Пусть $L_\rho(G)$ --- банахов функциональный $G$-модуль с абсолютно


непрерывной нормой $\rho$. Тогда


\begin{itemize}


\item[1)] $\cal L_\rho(G)\subset RUC(G);$


\item[2)] если $f\in\cal L_\rho(G)$, $g\in RUC(G)$ и $0\le g\le |f|$, то


$g\in\cal L_\rho(G);$


\item[3)] $\cal L_\rho(G)=RUC(G)$ в том и только том случае, когда


$L_\infty(G)\subset L_{\rho'}(G)$ {\em


(}или, эквивалентно, $L_\rho(G)\subset L_1(G));$


\item[4)] если $\cal L_\rho(G)\ne RUC(G)$, то


$\cal L_\rho(G)\subset\cal K(G)$.


\end{itemize}


\end{thms}





\begin{pf}


1),~2) Так как $f\circ F=f\ast\wt F$ для всех $f\in L_\rho(G)$,


$F\in L_{\rho'}(G)$, то $\cal L_\rho(G)$


совпадает с


замкнутым линейным подпространством в $L_\infty(G)$, порожденным


множеством $C_{00}(G)\ast L_{\rho'}(G)\sptilde$. Но


\begin{multline*}


C_{00}(G)\ast L_{\rho'}(G)\sptilde\subset L_{\rho'}(G)\sptilde


\cap\Cl\{C_{00}(G)\ast(C_{00}(G)\ast L_{\rho'}(G)\sptilde)\}\subset \\


\subset L_{\rho'}(G)\sptilde\cap RUC(G)\subset\Cl


\{C_{00}(G)\ast L_{\rho'}(G)\sptilde\},


\end{multline*}


следовательно,


\begin{equation}


\cal L_\rho(G)=\Cl\{L_{\rho'}(G)\sptilde\cap RUC(G)\}.


\end{equation}


Утверждения 1), 2) теперь очевидны.





3) В соответствии с (1) равенство $\cal L_\rho(G)=RUC(G)$ эквивалентно


тому, что $1_G\in L_{\rho'}(G)$. Но последнее условие означает, что


$L_\infty(G)\subset L_{\rho'}(G)$.





4) Так как $\cal L_\rho(G)\subsetneqq RUC(G)$, то


$1_G\notin\cal L_\rho(G)$. Следовательно,


существует


функционал $\Phi_0\in L_\infty(G)^*$ такой, что


$\Phi_0(1_G)=1$, $\Phi_0(f)=0$ для всех $f\in\cal L_\rho(G)$ и


$\|\Phi_0\|=1$.


Так как $\Phi_0$ является


топологически левоинвариантным на $\cal L_\rho(G)$ средним, то


$\cal L_\rho(G)$ --- аменабельный модуль.





Для любого топологически левоинвариантного на $\cal L_\rho(G)$ среднего


$\Phi$ и любой функции $f\in\cal L_\rho(G)$ имеем


$\la\varphi,|f|\circ\Phi\ra=\la\varphi^*\ast|f|,\Phi\ra=\Phi(|f|)$,


$\varphi\in P(G)$. Следовательно,


$\Phi(|f|)1_G=|f|\circ\Phi\in\cal L_\rho(G)$, откуда


$\Phi(|f|)=0$.


\end{pf}





\begin{cors}


Пусть $L_\rho(G)$ --- банахов функциональный $G$-модуль с абсолютно


непрерывной нормой $\rho$. Тогда следующие утверждения эквивалентны:


\begin{itemize}


\item[1)] модуль $L_\rho(G)$ аменабелен;


\item[2)] либо группа $G$ аменабельна, либо


$\cal L_\rho(G)\subset\cal K(G)$.


\end{itemize}


\end{cors}





\section{Модули со смешанной $L_p$-нормой}





Пусть $N_1$ --- замкнутая нормальная подгруппа в $G$, $f$ --- измеримая


функция на $G$. Обозначим через $X$ множество всех $t\in G$ таких, что


функция


$t_1\mapsto f(tt_1)$, $t_1\in N_1$, является


$\lambda_{N_1}$-измеримой. Тогда 1)~если


$t\in X$, то $tN_1\subset X$ и величина


$\int\limits_{N_1}|f(t't_1)|d\lambda_{N_1}(t_1)$ не зависит от выбора


$t'\in tN_1$; 2)~функция


$tN_1\mapsto\int\limits_{N_1}|f(tt_1)|d\lambda_{N_1}(t_1)$ является


$\lambda_{G/N_1}$-измеримой. Отсюда вытекает корректность следующего


определения.





Пусть $\Gamma:\{e\}=N_0\subset N_1\subset\dotsb\subset N_{n+1}=G$ ---


конечный возрастающий ряд замкнутых нормальных подгрупп,


$P=(p_0,\dotsc,p_n)\in[1,\infty)^{n+1}$. Произвольной измеримой на $G$


функции $f$ сопоставим величину


\begin{multline*}


\|f\|_P=\biggl(\int_{G/N_n}\dotsi\biggl(


\int_{N_2/N_1}\biggl(\int_{N_1}


|f(t_{n+1}\dotsb t_2t_1)|^{p_0}d\lambda_{N_1}(t_1)


\biggr)^{p_1/p_0}d\lambda_{N_2/N_1}(t_2N_1)\biggr)^{p_2/p_1}


\dotsb\\


\dotsb d\lambda_{G/N_n}(t_{n+1}N_n)\biggr)^{1/p_n}.


\end{multline*}


Если некоторые из величин $p_i$ равны $\infty$, в определение


$\|f\|_P$ необходимо внести соответствующие изменения.





Обозначим через $L_P(\Gamma)$ множество всех $\lambda_G$-измеримых


функций $f$,


удовлетворяющих условию $\|f\|_P<\infty$. Тогда


$L_P(\Gamma)$ является банаховым


функциональным $G$-модулем, причем


$$


\|f_s\|_P=\Delta_G(s)^{-1/p_0}\Delta_{G/N_1}


(sN_1)^{1/p_0-1/p_1}\dotsb\Delta_{G/N_n}


(sN_n)^{1/p_{n-1}-1/p_n}\|f\|_P.


$$


Общие свойства таких модулей изучались в [10].





В случае $P\in[1,\infty)^{n+1}$ справедливы следующие утверждения:


1)~норма $\|\,\cdot\,\|_P$ абсолютно


непрерывна ([11], теорема~3.1);


2)~$L_P(\Gamma)^*=L_{P'}(\Gamma)$, где


$P'=(p'_0,\dotsc,p'_n)$, $p^{-1}_i+p_i^{\prime-1}=1$ ([11], теорема~2.2).





Введем обозначение $\cal L_P(\Gamma)=\cal L_{L_P(\Gamma)}(G)$.





\begin{thms}


Пусть $P=[1,\infty)^{n+1}$, тогда следующие условия эквивалентны{\em :}


\begin{itemize}


\item[1)] модуль $L_P(\Gamma)$ аменабелен{\em ;}


\item[2)] либо группа $G$ аменабельна, либо для некоторого $i$


фактор-группа $N_{i+1}/N_i$ не является компактной и $p_i>1$.


\end{itemize}


\end{thms}





\begin{pf}


В соответствии с теоремой 3.3.\;


$\cal L_P(\Gamma)\subset\cal K(G)$ в том и только том случае, когда


$1_G\notin L_{P'}(\Gamma)$. Последнее условие эквивалентно тому, что для


некоторого $i$ фактор-группа


$N_{i+1}/N_i$ не является компактной и $p_i>1$. Осталось применить


следствие 3.1.


\end{pf}





При доказательстве теоремы 4.1 было отмечено, что


$\cal L_P(\Gamma)\subset\cal K(G)$ в том и только том


случае, когда для некоторого $i$ фактор-группа $N_{i+1}/N_i$ не


компактна и $p_i>1$.


Этот факт можно использовать для изучения ядра $\cal K(G)$. Ограничимся


наиболее простым случаем дискретной группы $G$.





Пусть $f\in L_\infty(G)$. Определим функции


$f^{(i)}\in L_\infty(N_{i+1}/N_i)$, $i=0,\dotsc,n$, равенством


$f^{(i)}(tN_i)=\|f(t\,\cdot\,)\|_{L_\infty(N_i)}$, $t\in N_{i+1}$. Для


произвольного


$\Omega\subset\{0,\dotsc,n\}$ обозначим через $C_{00}(\Gamma,\Omega)$


множество


функций $f\in L_\infty(G)$ таких, что


$\sup\limits_{s\in G}\card\supp({}_sf)^{(i)}<\infty$, $i\in \Omega$.


Пусть $C_0(\Gamma,\Omega)$ --- замыкание множества


$C_{00}(\Gamma,\Omega)$ в $L_\infty(G)$.





\begin{thms}


Пусть $G$ --- дискретная группа, $P\in[1,\infty)^{n+1}$ и


$\Omega=\{i:p_i>1\}$. Тогда $\cal L_P(\Gamma)=C_0(\Gamma,\Omega)$.


\end{thms}





\begin{pf}


Так как $C_{00}(\Gamma,\Omega)\subset L_{P'}(\Gamma)\subset


C_0(\Gamma,\Omega)$ и $C_{00}(\Gamma,\Omega)=


C_{00}(\Gamma,\Omega)\sptilde$,


то замыкание множества $L_{P'}(\Gamma)\sptilde$ в $L_\infty(\Gamma)$


совпадает с $C_0(\Gamma,\Omega)$. Осталось применить равенство (1).


\end{pf}





\begin{cors}


Пусть $G$ --- дискретная группа, и подмножество $\Omega$ таково, что


хотя бы для


одного $i\in\Omega$ фактор-группа $N_{i+1}/N_i$ бесконечна. Тогда


$C_0(\Gamma,\Omega)\subset\cal K(G)$.


\end{cors}





\begin{pf}


Выберем $P=(p_0,\dotsc,p_n)$ таким образом, чтобы $p_i>1$, $i\in\Omega$,


и $p_i=1$, $i\in\{0,\dotsc,n\}\setminus\Omega$. Тогда


$\cal L_P(\Gamma)\subset\cal K(G)$. Осталось


применить теорему 4.2.


\end{pf}





Предыдущее утверждение указывает на существование связи между ядром


компоненты аменабельности и решеткой замкнутых нормальных подгрупп.


\medskip





{\it Проблема}. Совпадает ли ядро $\cal K(G)$ с наименьшим замкнутым


подпространством в $L_\infty(G)$, содержащим все подмножества


$C_0(\Gamma,\Omega)$ указанного вида?








\begin{thebibliography}{99}





\bibitem{M1}


Мясников А.Г. {\em Аменабельные банаховы $L_1(G)$-модули, инвариантные


средние и регулярность в смысле Аренса} // Изв. вузов. Математика. --


1993. -- \No\,2. -- С.\,72--80.


\bibitem{Ros}


Rosenblatt J., Yang Z. {\em Functions with a unique mean value} // Math.


J. Illinois. -- 1990. -- V.\,34. -- \No\,4. -- P.\,744--764.


\bibitem{Mi1}


Miao T. {\em Amenability of locally compact groups and subspaces of


$L^\infty(G)$} // Proc. Amer. Math. Soc. -- 1991. -- V.\,111. -- \No\,4.


-- P.\,1075--1084.


\bibitem{Mi2}


Miao T. {\em The existence of left averaging functions that are not right


averaging} // Proc. Amer. Math. Soc. -- 1992. -- V.\,115. -- \No\,1.


-- P.\,121--123.


\bibitem{Em}


Emerson W.R. {\em Characterizations of amenable groups} // Trans. Amer.


Math. Soc. -- 1978. -- V.\,241. -- P.\,183--194.


\bibitem{Pi}


Pier J.-P. {\em Amenable locally compact groups}. -- New~York: Wiley,


1984. -- 418~p.


\bibitem{Pat}


Paterson A.L.T. {\em Amenability}. -- Providence, R.\,I.: AMS, 1988. --


452~p.


\bibitem{Ul}


Ulger A. {\em Continuity of weakly almost periodic functionals on


$L^1(G)$} // Quart J. Math. -- 1986. -- V.\,37. -- \No\,148. --


P.\,495--497.


\bibitem{M2}


Мясников А.Г. {\em Об абсолютной непрерывности норм в банаховых идеальных


пространствах со сдвигами} // Изв. вузов. Математика. -- 1990. -- \No\,3.


-- С.\,78--80.


\bibitem{M3}


Мясников А.Г., Черных С.И. {\em О пространствах $L_P$ со смешанной нормой


на локально компактных группах} // Функц. анализ и его прилож. -- 1985.


-- Т.\,19. -- \No\,3. -- С.\,73--74.


\bibitem{M4}


Мясников А.Г. {\em О смешанных нормах и абсолютно непрерывных нормах} //


1988. -- 24 с. -- Деп. в ВИНИТИ 02.06.88, \No\,4830-B88.








\end{thebibliography}





\vskip 2cm





\post{\it Московский государственный}{\it Поступила}


\post{\it строительный университет}{27.01.1995}





\end{document}


